% UNCOMPILED MATHEMATICAL PROOF TEXT. No compilation or verification was run.
% Intended as an insertion into the manuscript, using amsmath/amssymb/amsthm.
% Source scope: publication_package/proofs/quantum_results.tex, Section 13.
% Numbered local lemmas deliberately expose analytic dependencies for translation.

\section{Expanded analytic proofs for the modular response calculus}
\label{urgent:response}

\paragraph{Status and conventions.}
This text supplies expanded mathematical proofs, not a claim of
proof-assistant verification. It has not been compiled or computationally
checked. Matrices are finite-dimensional complex matrices, their norm is
the operator norm, and all matrix integrals can be read entry by entry.
For a self-adjoint matrix $K$, write $D=\operatorname{ad}_K$, so that
$D(X)=KX-XK$. Our Fourier convention is
$\widehat p(x)=\int_{\mathbb R}p(t)e^{itx}\,dt$.
All paths below are differentiable at the parameter under consideration;
continuously differentiable paths suffice throughout. No smooth choice of
eigenvectors along the path is assumed.

\subsection{Differentiation before spectral functional calculus}

\paragraph{Lemma R1 (the logarithm is differentiable).}
Let $X=X^*>0$, let $m>0$ satisfy $X\succeq mI$, and let $E=E^*$ with
$\|E\|<m/2$. Then
\[
 \log(X+E)-\log X
 =\int_0^\infty (X+sI)^{-1}E(X+E+sI)^{-1}\,ds,
\]
and
\[
 D\log_X(E)=\int_0^\infty(X+sI)^{-1}E(X+sI)^{-1}\,ds,
 \qquad
 \|\log(X+E)-\log X-D\log_X(E)\|
 \leq \frac{\|E\|^2}{m^2}.
\]
In particular, a differentiable positive matrix path $\rho$ has
$K=\log\rho$ differentiable, with
\[
 L:=\dot K=\int_0^\infty(\rho+sI)^{-1}\dot\rho
                                      (\rho+sI)^{-1}\,ds=L^*.
\]

\paragraph{Proof.}
The scalar identity, for $a>0$,
\[
 \log a=\int_0^\infty\left(\frac1{1+s}-\frac1{a+s}\right)ds
\]
follows by integrating to a finite endpoint and taking its limit.
Diagonalization proves the corresponding matrix identity. Set
$R_s=(X+sI)^{-1}$ and $R'_s=(X+E+sI)^{-1}$. The resolvent identity
$R'_s-R_s=-R_sER'_s$ proves the first displayed difference formula.
Its integrand is bounded by
$\|E\|/[(m+s)(m/2+s)]$, an integrable function.
Subtracting the candidate derivative gives exactly
\[
 \int_0^\infty R_sE(R'_s-R_s)\,ds
 =-\int_0^\infty R_sER_sER'_s\,ds.
\]
Because $\|R_s\|\leq(m+s)^{-1}$ and
$\|R'_s\|\leq2(m+s)^{-1}$, the norm of this remainder is at most
\[
 2\|E\|^2\int_0^\infty(m+s)^{-3}\,ds=\|E\|^2/m^2.
\]
The proposed derivative is bounded and linear, with norm at most $1/m$.
Apply this remainder bound to
$E=\rho(\lambda+h)-\rho(\lambda)=h\dot\rho+o(h)$.
Every integrand defining $L$ is self-adjoint. This also establishes the
existence of $\dot K$, rather than assuming it as an input.\hfill$\square$

\paragraph{Lemma R2 (exponential derivative and Duhamel formula).}
For arbitrary complex matrices $K,L$,
\[
 D\exp_K(L)
 =\sum_{n=1}^\infty\frac1{n!}\sum_{j=0}^{n-1}K^jLK^{n-1-j}
 =\int_0^1e^{(1-u)K}Le^{uK}\,du.
\]
Consequently, for $r\geq0$ fixed and $\rho=e^K>0$,
\begin{equation}
 (\partial_\lambda\rho^r)\rho^{-r}
 =\int_0^r e^{sK}\dot K e^{-sK}\,ds.
 \label{urgent:duhamel}
\end{equation}

\paragraph{Proof.}
The derivative of $K\mapsto K^n$ follows by the finite product rule;
it is the inner sum in the statement. On every norm ball
$\|K\|\leq R$, its norm is at most $nR^{n-1}\|L\|$.
The series of these bounds divided by $n!$ is $e^R\|L\|$.
The usual differentiation theorem for uniformly convergent derivative
series therefore applies to the exponential power series. This theorem
may also be applied entry by entry in the real and imaginary coordinates
of the matrix space.
For the integral expression, multiply the two absolutely uniformly
convergent exponential series on $0\leq u\leq1$ and integrate.
The elementary integer beta integral is
\[
 \int_0^1(1-u)^a u^b\,du=\frac{a!b!}{(a+b+1)!}
 \quad(a,b\in\mathbb Z_{\geq0});
\]
integration by parts in $b$, followed by the case $b=0$, proves it.
The coefficient of $K^aLK^b$ in the integral is thus
$1/(a+b+1)!$, giving the same series as the derivative.
Apply the derivative formula to $rK$, in direction $r\dot K$, and
multiply on the right by $e^{-rK}$. The change of variable
$s=r(1-u)$ proves \eqref{urgent:duhamel}; for $r=0$ both sides are zero.
\hfill$\square$

\paragraph{Lemma R3 (actual response and its spectral multipliers).}
Put $Q_r=(\partial_\lambda\rho^r)\rho^{-r}$ and
$H_r=(Q_r+Q_r^*)/2$, $A_r=(Q_r-Q_r^*)/2$. For real $x$, define
\[
 F_r(x)=\int_0^r e^{sx}\,ds,
 \quad S_r(x)=\begin{cases}\sinh(rx)/x&x\ne0,\\r&x=0,\end{cases}
 \quad T_r(x)=\begin{cases}(\cosh(rx)-1)/x&x\ne0,\\0&x=0.\end{cases}
\]
Then $Q_r=F_r(D)L$, $H_r=S_r(D)L$, and $A_r=T_r(D)L$.
Writing $H=H_1$, $A=A_1$, one obtains, for $0\leq r\leq1$,
\[
 H_r=\psi_r(D)H,\qquad
 A_r=\chi_r(D)A,
\]
where
\[
 \psi_r(x)=\frac{\sinh(rx)}{\sinh x},\quad\psi_r(0)=r,
 \qquad
 \chi_r(x)=\left(\frac{\sinh(rx/2)}{\sinh(x/2)}\right)^2,
 \quad\chi_r(0)=r^2.
\]
Moreover $L=f(D)H$, where $f(x)=x/\sinh x$ and $f(0)=1$, and
$Q_r/r\to L$ in norm as $r\downarrow0$.

\paragraph{Proof.}
At the single parameter value being differentiated, diagonalize $K$ with
real eigenvalues $\kappa_i$. Then $D$ acts on the matrix unit $E_{ij}$
by multiplication by $x=\kappa_i-\kappa_j$. Lemma R2 gives
$(Q_r)_{ij}=F_r(x)L_{ij}$. Since $L=L^*$ and $F_r$ is real on the
real line, $(Q_r^*)_{ij}=F_r(-x)L_{ij}$. Their half-sum and
half-difference are $S_r(x)L_{ij}$ and $T_r(x)L_{ij}$.
For $x\ne0$ the asserted ratios follow from
$\cosh(rx)-1=2\sinh^2(rx/2)$; at $x=0$ one checks the displayed
continuous values, with $T_r(0)=T_1(0)=0$.
The function $S_1$ is strictly positive on the real line, so its inverse
is $f$. Finally $F_r(x)/r\to1$ for every $x$. There are only finitely
many spectral differences, so the multiplier convergence is norm
convergence of matrix maps. This identifies the spectral expression with
the actual path derivative and does not presuppose that identification.
\hfill$\square$

\subsection{The scalar Fourier transform, including contour limits}

\paragraph{Lemma F1 (hyperbolic contour integral).}
For $0<\alpha<\pi$ and $k\in\mathbb R$,
\begin{equation}
 \int_{\mathbb R}\frac{e^{iku}}{\cosh u+\cos\alpha}\,du
 =\frac{2\pi}{\sin\alpha}\frac{\sinh(\alpha k)}{\sinh(\pi k)},
 \label{urgent:contour}
\end{equation}
where the quotient at $k=0$ means $\alpha/\pi$.

\paragraph{Proof.}
First take $k>0$ and integrate
$g(z)=e^{ikz}/(\cosh z+\cos\alpha)$ around the positively oriented
rectangle with vertices $-R,R,R+2\pi i,-R+2\pi i$.
The only poles inside are
$z_-=i(\pi-\alpha)$ and $z_+=i(\pi+\alpha)$.
Indeed the substitution $w=e^z$ reduces the denominator equation to
$w^2+2\cos\alpha\,w+1=0$, whose two roots have unit modulus;
the indicated arguments give exactly the roots in the open strip.
Both poles are simple, with derivatives
$\sinh z_-=i\sin\alpha$ and $\sinh z_+=-i\sin\alpha$.
Their residue sum is
\[
 \frac{e^{-k(\pi-\alpha)}-e^{-k(\pi+\alpha)}}{i\sin\alpha}
 =\frac{2e^{-\pi k}\sinh(\alpha k)}{i\sin\alpha}.
\]
The bottom integral is $I_R=\int_{-R}^R g(u)\,du$.
Periodicity of $\cosh$ makes the top integral
$-e^{-2\pi k}I_R$. On either vertical side,
\[
 |\cosh(\pm R+iy)|^2=\sinh^2R+\cos^2y,
 \qquad |e^{ik(\pm R+iy)}|=e^{-ky}\leq1.
\]
Thus, for sufficiently large $R$, the absolute value of each side
integral is at most $2\pi/(\sinh R-1)$, which tends to zero.
On the real axis the integrand is absolutely integrable: its denominator
is positive and is bounded below by a positive multiple of $e^{|u|}$
outside a compact interval. The residue theorem and $R\to\infty$ give
\[
 (1-e^{-2\pi k})I
 =\frac{4\pi e^{-\pi k}\sinh(\alpha k)}{\sin\alpha}.
\]
Divide using $1-e^{-2\pi k}=2e^{-\pi k}\sinh(\pi k)$ to obtain
\eqref{urgent:contour}. The denominator on the real axis is even, so the
sine part of the integral vanishes and $I(-k)=I(k)$.
Dominated convergence gives continuity at $k=0$, while
$\sinh(\alpha k)/\sinh(\pi k)\to\alpha/\pi$.
This proves all real cases, including the zero-frequency integral
$I(0)=2\alpha/\sin\alpha$.\hfill$\square$

\paragraph{Lemma F2 (positive fractional-power kernels).}
For $0<r<1$ define
\[
 p_r(t)=\frac{\sin(\pi r)}{2(\cosh(\pi t)+\cos(\pi r))}.
\]
Then $p_r\geq0$, $p_r\in L^1(\mathbb R)$, and
\[
 \widehat p_r(x)=\psi_r(x),\qquad \int p_r=r.
\]
The density $\widetilde p_r(t)=2p_r(2t)$ has transform
$\psi_r(x/2)$, and $\widetilde p_r*\widetilde p_r$ has transform
$\chi_r(x)$ and mass $r^2$.
Furthermore
\[
 \frac{p_r}{r}\longrightarrow
 p_0(t):=\frac\pi4\operatorname{sech}^2(\pi t/2)
 \quad\hbox{in }L^1,\qquad
 \widehat p_0(x)=f(x),\quad\int p_0=1.
\]

\paragraph{Proof.}
Apply Lemma F1 with $\alpha=\pi r$, $u=\pi t$, $k=x/\pi$,
including its $k=0$ conclusion. Positivity is immediate since
$\sin(\pi r)>0$ and $\cosh(\pi t)+\cos(\pi r)>0$.
The rescaling formula follows by $v=2t$; the convolution formula and
mass statement follow by Fubini, justified by the product of the two
$L^1$ norms. For $0<r\leq1/2$,
\[
 0\leq p_r(t)/r\leq\pi/[2\cosh(\pi t)],
\]
because $\sin(\pi r)/r\leq\pi$ and $\cos(\pi r)\geq0$.
This bound is integrable, and the pointwise limit is
$\pi/[2(\cosh(\pi t)+1)]=p_0(t)$.
Dominated convergence gives the stated $L^1$ convergence and mass.
Fourier transforms converge under $L^1$ convergence; the scalar limit
$\psi_r(x)/r\to f(x)$ identifies the transform.\hfill$\square$

\paragraph{Corollary F3 (complete positivity and exact norms).}
For self-adjoint $K$, the maps $\psi_r(D)$ and $\chi_r(D)$ are
completely positive and have completely bounded operator norms $r$ and
$r^2$, respectively. The map $f(D)$ is unital completely positive.

\paragraph{Proof.}
In the eigenbasis of $K$, Fourier inversion in the sense of Lemma F2
(namely its forward integral identity) gives exactly
\[
 \psi_r(D)(X)=\int p_r(t)e^{itK}Xe^{-itK}\,dt.
\]
For every finite ancilla the same formula holds with
$e^{itK}\otimes I$. A nonnegative integral of positive matrices is
positive, proving positivity at every amplification. The operator-norm
bound at each amplification is the mass of the density, by the triangle
inequality and unitary invariance. Evaluation at the identity attains
that mass. Apply the same argument to the convolution density and $p_0$.
At $r=0$ the maps are zero; at $r=1$ both are the identity. Thus no
endpoint density formula is needed.\hfill$\square$

\subsection{Deriving the coth series instead of assuming it}

\paragraph{Lemma C1 (an elementary partial-fraction identity).}
For every real $x\ne0$,
\begin{equation}
 x\coth(x/2)=2+4\sum_{n=1}^\infty
                         \frac{x^2}{x^2+(2\pi n)^2}.
 \label{urgent:coth-series}
\end{equation}
The right-hand series and its first derivative converge uniformly on
every bounded real interval after assigning the value $2$ at zero.

\paragraph{Proof.}
Fix $a>0$. Regard $t\mapsto\cosh(at)$ on $[-\pi,\pi]$ as a
continuous $2\pi$-periodic function; its endpoint values agree.
Its sine coefficients vanish by evenness. Direct integration gives
\[
 \frac1{2\pi}\int_{-\pi}^{\pi}\cosh(at)\,dt
       =\frac{\sinh(a\pi)}{\pi a},
\]
and, for $n\geq1$,
\[
 \frac1\pi\int_{-\pi}^{\pi}\cosh(at)\cos(nt)\,dt
       =\frac{2a\sinh(a\pi)(-1)^n}{\pi(a^2+n^2)}.
\]
For the second equality one differentiates the explicit antiderivative
\[
 \frac{a\sinh(at)\cos(nt)+n\cosh(at)\sin(nt)}{a^2+n^2}.
\]
The resulting cosine series is absolutely uniformly convergent, since
its coefficients are bounded by a constant times $n^{-2}$.
It equals the continuous periodic function: for example, the uniform
Fej\'er convergence theorem identifies the Ces\`aro sums of its Fourier
series with that function, while absolute uniform convergence identifies
the same Ces\`aro limit with the displayed series.
Evaluating at $t=\pi$ therefore gives
\[
 \cosh(a\pi)=\frac{\sinh(a\pi)}{\pi a}
       +\frac{2a\sinh(a\pi)}\pi\sum_{n\geq1}\frac1{a^2+n^2}.
\]
Divide by $\sinh(a\pi)$ and put $x=2\pi a$. This is
\eqref{urgent:coth-series} for $x>0$; both sides are even, giving $x<0$.
Write $b_n=2\pi n$. On $|x|\leq R$, the summand
$4x^2/(x^2+b_n^2)$ is bounded by $4R^2/b_n^2$, and its derivative
$8xb_n^2/(x^2+b_n^2)^2$ by $8R/b_n^2$.
The Weierstrass bounds prove both asserted uniform convergences and
justify termwise differentiation.\hfill$\square$

\paragraph{Lemma C2 (positive series for the rotation multiplier).}
Define
\[
 h(x)=\frac{1-x/\sinh x}{x\tanh(x/2)}\quad(x\ne0),
 \qquad h(0)=\frac13.
\]
Then, for all real $x$,
\begin{equation}
 h(x)=8\sum_{n=1}^\infty
                    \frac{(2\pi n)^2}{(x^2+(2\pi n)^2)^2}.
 \label{urgent:h-series}
\end{equation}

\paragraph{Proof.}
For $x\ne0$, elementary hyperbolic identities give
\[
 \frac{1-x/\sinh x}{x\tanh(x/2)}
 =\frac{\coth(x/2)}x-\frac12\operatorname{csch}^2(x/2)
 =\frac1x\frac d{dx}\{x\coth(x/2)\}.
\]
Lemma C1 permits differentiating the series, giving
\eqref{urgent:h-series} for nonzero $x$. Each summand is bounded
in absolute value by $8/(2\pi n)^2$, so the series is continuous at zero.
The ordinary Taylor expansions of $\sinh$ and $\cosh$, obtained from
their exponential series with remainders, give
$x\coth(x/2)=2+x^2/6+O(x^4)$ and
$\frac d{dx}\{x\coth(x/2)\}=x/3+O(x^3)$.
Thus its derivative divided by $x$ tends to $1/3$, establishing the
zero value too. This also proves
$\sum_{n\geq1}8/(2\pi n)^2=1/3$ without importing a Basel-sum value.
\hfill$\square$

\paragraph{Lemma C3 (explicit nonnegative Fourier density).}
For $t\ne0$ let
\[
 k(t)=\sum_{n=1}^\infty
       \left(\frac1{\pi n}+2|t|\right)e^{-2\pi n|t|};
\]
assign any finite value, for example zero, at $t=0$. Then $k\geq0$
almost everywhere, $k\in L^1$, $\int k=1/3$, and
\[
 \widehat k(x)=h(x).
\]
For $t\ne0$ it also has the closed form
\[
 k(t)=-\frac1\pi\log(1-e^{-2\pi|t|})
              +\frac{2|t|}{e^{2\pi|t|}-1}.
\]
For every real $v<2\pi$ its exponentially weighted mass is finite;
for $0\leq v<2\pi$ it equals
\begin{equation}
 \int_{\mathbb R}e^{v|t|}k(t)\,dt
 =2\sum_{n\geq1}\left\{
 \frac1{\pi n(2\pi n-v)}+
 \frac2{(2\pi n-v)^2}\right\}.
 \label{urgent:weighted-mass}
\end{equation}

\paragraph{Proof.}
For $a>0$, elementary exponential integration on each half-line yields
\[
 \int_{\mathbb R}e^{-a|t|}e^{itx}dt=\frac{2a}{a^2+x^2},
\qquad
 \int_{\mathbb R}|t|e^{-a|t|}e^{itx}dt
       =\frac{2(a^2-x^2)}{(a^2+x^2)^2}.
\]
For completeness, integrate $e^{-(a-ix)t}$ and $t e^{-(a-ix)t}$
on $[0,R]$, let $R\to\infty$, and add their complex conjugates;
the boundary terms vanish since $a>0$. It follows that
\[
 \int_{\mathbb R}(2/a+2|t|)e^{-a|t|}e^{itx}dt
     =\frac{8a^2}{(a^2+x^2)^2}.
\]
The mass of this nonnegative function is $8/a^2$.
Set $a=2\pi n$ and sum. Tonelli and Lemma C2 give total mass $1/3$.
The summed integrable masses justify exchanging the complex Fourier
integral with the series, proving the transform identity.
For the closed form use $\sum_{n\geq1}q^n=q/(1-q)$ and
$\sum_{n\geq1}q^n/n=-\log(1-q)$ for $0<q<1$; the latter follows
by integrating the geometric series on a compact subinterval of $[0,1)$.
The value at the single point zero does not affect any integral.
Finally Tonelli and direct exponential integration give
\eqref{urgent:weighted-mass}. Since
$2\pi n-v\geq(2\pi-v)n$ for $n\geq1$ and $0\leq v<2\pi$,
its terms are bounded by a fixed constant times $n^{-2}$.
Negative $v$ only decrease the integrand.\hfill$\square$

\subsection{KMS compatibility and the rotation identity}

\paragraph{Lemma K1 (compatibility of the actual derivative).}
For $Q=\dot\rho\rho^{-1}=H+A$, with $H=H^*$ and $A^*=-A$,
\[
 Q^*=e^{-D}Q,\qquad A=\tanh(D/2)H.
\]

\paragraph{Proof.}
Because $\rho$ and $\dot\rho$ are self-adjoint,
$Q^*=\rho^{-1}\dot\rho=\rho^{-1}Q\rho=e^{-D}Q$.
Substitute $Q=H+A$ and $Q^*=H-A$ to obtain
$(I-e^{-D})H=(I+e^{-D})A$.
In an eigenbasis of $K$, every eigenvalue of $I+e^{-D}$ is
$1+e^{-(\kappa_i-\kappa_j)}>0$, so it is invertible.
The quotient is $(1-e^{-x})/(1+e^{-x})=\tanh(x/2)$, proving the claim.
\hfill$\square$

\paragraph{Proposition K2 (positive rotation generator, with all inputs proved).}
Put $J=-iA=J^*$ and $C=h(D)J$. Then
\[
 C=C^*,\qquad \|C\|\leq\tfrac13\|J\|,
 \qquad \dot K=H+i[C,K].
\]
The map $h(D)$ is completely positive, has completely bounded norm
$1/3$, and admits the exponentially weighted Fourier mass bounds of
Lemma C3. If a differentiable unitary path solves $\dot U=iCU$, then
\[
 \frac d{d\lambda}(U^*KU)=U^*HU.
\]

\paragraph{Proof.}
Lemma C3 gives
$h(D)(X)=\int k(t)e^{itK}Xe^{-itK}\,dt$.
The argument of Corollary F3 proves complete positivity, preservation
of adjoints, and the exact norm $\int k=1/3$.
By Lemma K1, $J=-i\tanh(D/2)H$. The scalar identity, including its
continuous zero value, is
$xh(x)\tanh(x/2)=1-f(x)$.
Consequently
\[
 i[C,K]=-iD h(D)J
       =-D h(D)\tanh(D/2)H
       =f(D)H-H=\dot K-H,
\]
where the last equality uses Lemma R3. Finally the product rule and
$\dot U^*=-iU^*C$ give
\[
 \frac d{d\lambda}(U^*KU)
 =U^*\{\dot K-i[C,K]\}U=U^*HU.
\]
If $C$ is continuous, the standard finite-dimensional linear ODE theorem
constructs $U$ with $U(0)=I$; differentiating $U^*U$ shows it remains
unitary. The response and rotation conclusions themselves do not require
that auxiliary ODE construction.\hfill$\square$

\paragraph{Composition and square roots.}
For $r,s\in[0,1]$, direct cancellation of hyperbolic ratios gives
\[
 \psi_s(rx)\psi_r(x)=\psi_{rs}(x),\qquad
 \chi_s(rx)\chi_r(x)=\chi_{rs}(x).
\]
Check $x=0$ and $r=0$ from the prescribed continuous values.
The identities $\sinh x=2\sinh(x/2)\cosh(x/2)$ then give
$\psi_{1/2}(x)=\tfrac12\operatorname{sech}(x/2)$ and
$\chi_{1/2}(x)=\tfrac14\operatorname{sech}^2(x/4)$.
Thus the expanded proofs also recover the composition, endpoint,
ancilla, and square-root statements in the report.
